% packages_documentation.tex
% Official sources + reading guide for every external package/tool used in the
% Conway order-7 non-existence artifact (certified proof that 7 does not divide
% |Aut(Gamma)| for any srg(99,14,1,2)).
%
% Compile with: pdflatex packages_documentation.tex

\documentclass[11pt]{article}

\usepackage[a4paper,margin=1in]{geometry}
\usepackage[hidelinks]{hyperref}

\newcommand{\tool}[1]{\subsection{#1}}
\newcommand{\site}[2]{\textbf{#1:} \url{#2}\\}

\title{Packages \& Tools Used in the\\ Conway Order-7 Non-Existence Artifact\\[4pt]
       \large Official Sources and Reading Guide}
\author{}
\date{}

\begin{document}
\maketitle

\noindent
This document lists every external package, binary, and format the artifact
depends on, with the \emph{official} website/repository for each and a pointer
to the key paper or documentation to read to understand how it works. The
mathematical claim is: every one of the 98{,}536 cubes is UNSAT (each backed by
a machine-checkable LRAT certificate) and the cubes exactly tile the search
space, hence no $\mathrm{srg}(99,14,1,2)$ admits an order-7 automorphism. The SAT
solver is \emph{untrusted}; trust rests only on the proof checkers and the
coverage arithmetic --- so the checkers below are the most important tools to
understand.

\section*{At a glance}
\noindent\begin{center}\small
\begin{tabular}{@{}p{0.22\textwidth} p{0.30\textwidth} p{0.40\textwidth}@{}}
\textbf{Tool / file} & \textbf{Role in this artifact} & \textbf{Official source} \\
\hline
CaDiCaL (\texttt{tools/cadical}) & SAT solver; refutes each cube, emits LRAT & \url{https://github.com/arminbiere/cadical} \\
PySAT (\texttt{python-sat}) & CNF encoder library (\texttt{conway\_o7.py}) & \url{https://pysathq.github.io} \\
lrat-check (\texttt{tools/lrat-check}) & LRAT checker (Heule) --- trusted base & \url{https://github.com/marijnheule/drat-trim} \\
lrat-trim (\texttt{tools/lrat-trim}) & LRAT checker/trimmer (Biere) & \url{https://github.com/arminbiere/lrat-trim} \\
cake\_lpr (planned) & Formally-verified LRAT checker & \url{https://github.com/tanyongkiam/cake_lpr} \\
DRAT/LRAT format & The proof format itself & \url{https://github.com/marijnheule/drat-trim} \\
Zstandard (\texttt{zstd}) & Compresses the 98{,}536 proofs \& the kit & \url{https://facebook.github.io/zstd/} \\
Slurm & HPC job scheduler (the cluster run) & \url{https://slurm.schedmd.com/} \\
Apptainer / Singularity & Container for the cluster / re-run & \url{https://apptainer.org/} \\
Python 3 & Runs encoder, coverage, aggregation & \url{https://www.python.org/} \\
\end{tabular}
\end{center}

\section{SAT solving and proof certification (the core)}

\tool{CaDiCaL --- the SAT solver}
The complete CDCL SAT solver by Armin Biere. In this artifact it refutes each
cube and, run with \texttt{-{}-factor=false}, emits a plain-RUP LRAT proof
(\texttt{solve\_cube.py --proof}). It is also invoked \emph{inside} the encoder
through PySAT's \texttt{Cadical195} interface. Remember: CaDiCaL is
\textbf{untrusted} here --- if it lies, the checkers below catch it.
\begin{description}
  \item[Official repository] \url{https://github.com/arminbiere/cadical}
  \item[Project page (Biere / FMV)] \url{http://fmv.jku.at/cadical/}
  \item[Read to understand it] The repository \texttt{README.md}, and the tool
        paper: A.\ Biere, T.\ Faller, K.\ Fazekas, M.\ Fleury, N.\ Froleyks,
        F.\ Pollitt, ``CaDiCaL 2.0'' (CAV 2024),
        \url{https://link.springer.com/chapter/10.1007/978-3-031-65627-9_7}
\end{description}

\tool{PySAT (\texttt{python-sat}) --- the CNF encoder library}
The Python toolkit used by \texttt{conway\_o7.py}/\texttt{export\_o7.py} to build
the DIMACS instance \texttt{o7.cnf}. It provides the cardinality encodings for
the degree and common-neighbour constraints (the project selects the
\emph{modulo totalizer}, \texttt{mtotalizer}, via
\texttt{CONWAY\_DEG\_ENC=mtotalizer}) and a uniform interface to CaDiCaL.
\begin{description}
  \item[Documentation] \url{https://pysathq.github.io}
  \item[Repository] \url{https://github.com/pysathq/pysat}
  \item[PyPI package] \url{https://pypi.org/project/python-sat/}
  \item[Cardinality encodings API] \url{https://pysathq.github.io/docs/html/api/card.html}
  \item[Read to understand it] A.\ Ignatiev, A.\ Morgado, J.\ Marques-Silva,
        ``PySAT: A Python Toolkit for Prototyping with SAT Oracles'' (SAT 2018) ---
        linked from the docs homepage. For the encodings themselves see the
        totalizer / modulo-totalizer references on the \texttt{pysat.card} page.
\end{description}

\tool{The DRAT / LRAT proof format}
LRAT (``Linear RAT'') is the machine-checkable UNSAT-proof format the whole
certificate is built on: each line is a resolution/RAT step with explicit
clause-id hints, so a checker verifies it in one linear pass without search.
This is the format \texttt{cadical} emits and every checker below consumes.
Understand this format \emph{first} --- it is the object of trust.
\begin{description}
  \item[DRAT-trim / lrat-check home] \url{https://github.com/marijnheule/drat-trim}
  \item[LRAT format paper] L.\ Cruz-Filipe, M.\ Heule, W.\ Hunt, M.\ Kaufmann,
        P.\ Schneider-Kamp, ``Efficient Certified RAT Verification'' (CADE 2017).
  \item[DRAT background] N.\ Wetzler, M.\ Heule, W.\ Hunt, ``DRAT-trim: Efficient
        Checking and Trimming Using Expressive Clausal Proofs'' (SAT 2014).
\end{description}

\tool{lrat-check (Heule) --- checker \#1}
The reference LRAT checker (\texttt{lrat-check.c}) that ships in the DRAT-trim
repository. It is what \texttt{verify\_final.sh} / \texttt{verify\_proofs.sh}
call to re-verify every certificate. Small, public, and replaceable --- part of
the trusted base.
\begin{description}
  \item[Source (in DRAT-trim repo)] \url{https://github.com/marijnheule/drat-trim/blob/master/lrat-check.c}
  \item[Read to understand it] The header comments in \texttt{lrat-check.c}, plus
        the LRAT format paper above.
\end{description}

\tool{lrat-trim (Biere) --- checker \#2 / trimmer}
Armin Biere's independent LRAT reader/checker/trimmer. Using it \emph{as well as}
lrat-check means two independently-written checkers agree on every proof
(defence against a bug in either one).
\begin{description}
  \item[Repository] \url{https://github.com/arminbiere/lrat-trim}
  \item[Read to understand it] The repository \texttt{README} and the top of
        \texttt{lrat-trim.c}; background paper: F.\ Pollitt, M.\ Fleury,
        A.\ Biere, ``Faster LRAT Checking Than Solving with CaDiCaL'' (SAT 2023),
        \url{https://drops.dagstuhl.de/entities/document/10.4230/LIPIcs.SAT.2023.21}
\end{description}

\tool{cake\_lpr (CakeML) --- checker \#3, formally verified (planned)}
The strongest check on the to-do list: a checker whose x86-64 machine code is
\emph{itself} proved correct in the CakeML/HOL4 ecosystem. Because all proofs
are plain-RUP LRAT, cake\_lpr can be dropped in as a third, formally-verified
verification pass.
\begin{description}
  \item[Repository] \url{https://github.com/tanyongkiam/cake_lpr}
  \item[CakeML project] \url{https://cakeml.org}
  \item[Read to understand it] Y.\ K.\ Tan, M.\ Heule, M.\ O.\ Myreen,
        ``cake\_lpr: Verified Propagation Redundancy Checking in CakeML''
        (TACAS 2021), \url{https://cakeml.org/tacas21.pdf}
\end{description}

\section{Methodology (not a package, but essential to understand)}

\tool{Cube-and-Conquer}
The parallelisation strategy behind \texttt{cubes.txt}, \texttt{deepen.py}, and
the per-cube UNSAT proofs: split (``cube'') the search space into many
sub-problems, then conquer each with a full CDCL solver. Understanding this
explains why there are 98{,}536 independent proofs and why coverage
(\texttt{check\_coverage.py}, exact mass $=1$) is a necessary second ingredient.
\begin{description}
  \item[Paper] M.\ Heule, O.\ Kullmann, S.\ Wieringa, A.\ Biere, ``Cube and
        Conquer: Guiding CDCL SAT Solvers by Lookaheads'' (HVC 2011).
  \item[Author hub (many related papers)] Marijn Heule, \url{https://www.cs.cmu.edu/~mheule/}
\end{description}

\tool{Lean 4 + mathlib --- theorem prover (formalization layer, in progress)}
The proof assistant hosting the project's formalization layer
(\texttt{Conway\_lean/} in this folder, Lean v4.30.0-rc2 + mathlib):
theorem T1 (the 98{,}536 cubes cover the space, replacing
\texttt{check\_coverage.py} in the trusted base) and, later, theorem T2 (the CNF
encoding is faithful: UNSAT $\Rightarrow$ no $\sigma$-invariant srg). cake\_lpr
keeps discharging the per-cube UNSAT hypotheses; Lean proves the reasoning
around them.
\begin{description}
  \item[Lean website] \url{https://lean-lang.org/}
  \item[mathlib (community math library)] \url{https://leanprover-community.github.io/}
  \item[Read to understand it] ``Theorem Proving in Lean 4'',
        \url{https://lean-lang.org/theorem_proving_in_lean4/}
\end{description}

\section{Compression and infrastructure}

\tool{Zstandard (zstd)}
The compressor for the 98{,}536 \texttt{*.lrat.zst} proofs and the
\texttt{referee\_kit.tar.zst} supplement; \texttt{verify\_final.sh} decompresses
on the fly.
\begin{description}
  \item[Website] \url{https://facebook.github.io/zstd/}
  \item[Repository] \url{https://github.com/facebook/zstd}
\end{description}

\tool{Slurm Workload Manager}
The HPC batch scheduler that ran the cube array on the cluster
(\texttt{o7\_conquer.slurm}, \texttt{SLURM\_TMPDIR}).
\begin{description}
  \item[Documentation] \url{https://slurm.schedmd.com/}
\end{description}

\tool{Apptainer / Singularity (and Docker)}
Containers appear in \emph{two separate} roles here. (1)~\textbf{Past --- proof
generation (finished).} Singularity was the container runtime on the HPC cluster
(\texttt{module load singularity}) that ran the cube-and-conquer solve; that
compute is done and is not repeated. (2)~\textbf{Planned --- reproducible
verification.} A small Docker/Apptainer image bundling PySAT + CaDiCaL + the
checkers, so a referee can re-run the \emph{verification} (not the search) with
one command (the \texttt{container/} folder). Note: Apptainer is simply the
renamed Singularity, and it can also run Docker images.
\begin{description}
  \item[Apptainer (current name of Singularity)] \url{https://apptainer.org/}
  \item[Docker (alternative)] \url{https://www.docker.com/}
\end{description}

\tool{Python 3}
Runs the encoder, the exact-arithmetic coverage checker, and the aggregation
scripts; the \texttt{refvenv} virtual environment holds PySAT.
\begin{description}
  \item[Website] \url{https://www.python.org/}
\end{description}

\end{document}
